%=======================================================================================
% MERGED/VERIFIED DATA CHAPTER -- 2026-08-10 : 12:33pm
%
% Base content: documentation/august_draft/3_Chapter_data/data_chapter_draft.tex, which already
% uses the same chapter/section layout as the dissertation.tex skeleton pasted into the planning
% session (\section{Forest research data}, \section{Data Cleaning}, \section{stand terrain and
% environmental data}, \section{Terrain data}, \section{Wind data}, \section{Climate, soil and
% other spatial data}) -- no structural merge was needed here, unlike Methodology.
% This pass: cross-checked every specific number against data_processing/clean_master_data.py
% (read directly, 2026-08-10) and documentation/progress_notes.md / experiment_log.md. Numbers
% that matched the code/notes exactly are left as-is. Numbers that could not be found in current
% artifacts are marked NEEDS CHECKING inline, not silently kept or silently dropped.
% Results-chapter content excluded, per this project's own rule (Data states facts, Results
% states findings) -- e.g. the HadUK-Grid Moran's I number stays here only because it is used to
% justify a data-resolution decision, not as an attribution finding.
%=======================================================================================

\chapter{Data}
\section{Forest research data}
Data supplied by Forest Research as a GeoPackage.
Study area around Aberfoyle.
British National Grid, EPSG:27700.
Repeated forest-plot polygons from 2002, 2006, 2008, 2012, 2021 and 2023.
Each record represents a plot in one survey year.
Plots are roughly all 20m by 20m (area field median exactly 400\,m\textsuperscript{2} across all cohort plots), with smaller irregular exceptions (down to around 250\,m\textsuperscript{2}) where a plot's sampling grid cell is clipped by a compartment or sub-compartment boundary.
Plot and compartment are different levels of the same forestry hierarchy, not the same unit measured two ways: a compartment (\texttt{cpmt}) is a Forest Research management unit, further divided into sub-compartments (\texttt{scpt}); a plot is the individual LiDAR/Area-Based-Analysis sampling grid cell that survey statistics were computed over, 
and many plots sit inside one compartment.
The full data contain 795,645 plot-year records and 230,359 plot identifiers.
Sitka spruce was identified using \texttt{spis == "SS"}.
The main target was the 95th-percentile LiDAR height (\texttt{elev\_percentile\_95th}, confirmed as \texttt{TARGET\_COLUMN} in \texttt{data\_processing/clean\_master\_data.py}) -- raw, no correction applied, distinct from the retired \texttt{Top\_Height99}/$\times1.1$ pathway.
Important original variables included planting year, canopy cover, thinning information, species, yield class and windthrow hazard class.
The height data are Airborne Laser Scanning (ALS), collected from aircraft, not drone-based LiDAR.

\textbf{NEEDS CHECKING:} the master export (\texttt{MASTER\_COLUMNS} in \texttt{clean\_master\_data.py}) also retains \texttt{LAI}, but no plausibility/sanity filter is applied to it at the cleaning stage (only \texttt{Age} and the height target are range-checked) -- if an earlier data-quality pass flagged an anomaly in this column for a specific survey year, that has not been re-verified against the current cleaned cohorts and should be checked before \texttt{LAI} is used anywhere downstream.

\begin{table}[h]
\centering
\begin{tabular}{p{2.2cm}p{3cm}p{3cm}p{3cm}p{3.5cm}}
\hline
Data family & Source & Resolution or structure & Temporal representation & Main purpose \\
\hline
Forest plots & Forest Research & Repeated plot polygons, roughly 20m by 20m & Six survey years, 2002--2023 & Prediction and growth calculation \\
Terrain & OS Terrain 50 & 50\,m grid & Static & Topography and exposure \\
Climate & HadUK-Grid & 1\,km grid (149 distinct values across 71,766 plots) & Annual or multi-year cohort-averaged summaries & Temperature, frost and rainfall conditions \\
Wind climate & Global Wind Atlas & Raster climatology, nominal 250\,m (1,816 distinct values across 71,766 plots) & Long-term static summaries & General wind exposure \\
Observed wind & MIDAS stations & Station observations & Survey-interval summaries & Temporal wind and storm conditions \\
Other environment & CHELSA, CEH, SoilGrids and spatial layers & Source-dependent, 250\,m--1\,km & Mainly static & Wider environmental screening \\
\hline
\end{tabular}
\end{table}
\textbf{Resolved (2026-08-10):} the earlier draft's "around 1,800 distinct values" figure for Global Wind Atlas is confirmed -- \texttt{notebooks/environmental\_data/figures/aux\_data\_resolution\_check\_results.csv} (the notebook's own saved numeric output, read directly, not the notebook's prose) records \texttt{n\_distinct\_values=1816} for "Global Wind Atlas (10m wind speed)" across the same 71,766-plot set used elsewhere in this chapter. The nominal $0.0025^\circ$ grid is \textbf{anisotropic} at Aberfoyle's latitude ($\approx$56.18\textdegree N) -- narrower east--west than north--south, since a degree of longitude covers less ground at higher latitude. \textbf{Recomputed directly from the cached raster} (\texttt{data/raw/environmental/gbr\_wind\_10m.tif}, \texttt{rasterio} transform, 2026-08-10): confirmed pixel size $0.0025^\circ\times0.0025^\circ$, giving \textbf{154.9\,m east--west $\times$ 278.3\,m north--south} at lat.\ 56.18\textdegree N. This matches the current resolution-check CSV's $\approx$155\,m/$\approx$278\,m figures; a superseded exploratory notebook's "$\approx$140--155\,m E--W $\times$ 250\,m N--S" note undercounted the north--south dimension (the flat "250m" marketing figure, not the latitude-corrected value) -- an example of why this pass trusts recomputed numbers over either notebook's stated prose.

\section{Data Cleaning}
The cleaning funnel builds two balanced cohorts from the same raw layer.
Confirmed directly against \texttt{clean\_master\_data.py}'s \texttt{clean\_cohort()} function (2026-08-10) -- the actual coded stage order is: (0) raw dataset $\rightarrow$ (1) restrict to the cohort's chosen survey years $\rightarrow$ (2) complete coverage across every one of those years $\rightarrow$ (3) only Sitka spruce (\texttt{spis == "SS"}) $\rightarrow$ (4) a valid planting year (\texttt{plyr} $\ge$ 1850 and strictly before the survey year) $\rightarrow$ (5) plausible age $\rightarrow$ (6) plausible height.
At every stage from (3) onward, a whole \emph{plot} is dropped only if any one of its survey-year rows fails that stage's condition -- balanced trajectories are preserved, not partial ones.
The plausibility rules themselves, confirmed in code: \texttt{MIN\_PLYR = 1850}; \texttt{Age} (survey year minus planting year) must fall in \texttt{[0, 200]}; height (the target column, \texttt{elev\_percentile\_95th}) must be $>0$\,m and $\le 60$\,m -- a generous ceiling for mature Sitka spruce, and the code's own comment notes these three bounds are wide sanity checks, not domain-cited limits, chosen so they should never bind on real data at this site.

\begin{table}[h]
\centering
\begin{tabular}{lrr}
\hline
Stage or characteristic & Four-survey cohort & Six-survey cohort \\
\hline
Complete plots before species filtering & 112,782 & 18,781 \\
Sitka spruce plots & 72,395 & 13,897 \\
Final plots & 71,766 & 13,897 \\
Final plot-year records & 287,064 & 83,382 \\
Survey years & 2008, 2012, 2021, 2023 & 2002, 2006, 2008, 2012, 2021, 2023 \\
Observations per plot & 4 & 6 \\
Survey intervals & 4, 9 and 2 years & 4, 2, 4, 9 and 2 years \\
Role & Main analysis & Sensitivity analysis \\
\hline
\end{tabular}
\end{table}
\textbf{Confirmed (2026-08-10):} re-ran \texttt{python -m data\_processing.clean\_master\_data} directly against the raw GeoPackage (795,645 rows, 230,359 plots read) -- every row-count value in the table above matches the fresh run exactly, including the intermediate "1. Cohort years" stage not previously shown (4survey: 735,219 rows/230,345 plots after restricting to 2008/2012/2021/2023; 6survey: unchanged at 795,645/230,359, since six-survey uses every year in the raw data). One sharper fact this re-run surfaces that the draft's own text did not state precisely: stages 5--6 (plausible age, plausible height) remove \textbf{zero} additional rows in \emph{either} cohort, and stage 4 (valid planting year) removes rows only in the \textbf{four}-survey cohort (289,580$\to$287,064 rows, 72,395$\to$71,766 plots) -- the six-survey cohort loses \emph{nothing at all} from stage 3 onward (83,382 rows/13,897 plots unchanged through stages 3--6). \texttt{thinning\_status=`unknown\_timing'} (corrupt \texttt{last\_thinn}) count: confirmed 0 rows, both cohorts.

The four-survey cohort has more plots, so greater spatial coverage.
The six-survey cohort has two more observations per plot but retained fewer than one-fifth as many plots.
The four-survey cohort was therefore used for the main analysis; the six-survey cohort tested whether more years changed the conclusions.

\section{stand terrain and environmental data}
\subsection{stand and management}
Canopy cover represented the observed canopy condition.
Recorded thinning information was used to derive thinning status, time since thinning and recent-thinning indicators -- confirmed in \texttt{clean\_master\_data.py}'s \texttt{add\_export\_metrics()}: \texttt{Thin} is 1 only if \texttt{last\_thinn} is a valid, already-occurred year ($0 < \text{last\_thinn} \le$ survey year); \texttt{time\_since\_thinning}, 
\texttt{time\_since\_thinning\_missing}, \texttt{recent\_thinning\_5yr}/\texttt{10yr}, and a five-level \texttt{thinning\_status} (\texttt{never\_thinned}, \texttt{recent\_0\_5yr}, \texttt{mid\_6\_10yr}, \texttt{old\_10plus\_yr}, \texttt{not\_yet\_thinned\_at\_survey}, plus an \texttt{unknown\_timing} catch-all for a corrupt \texttt{last\_thinn} value) 
are all derived per survey row, not per plot -- so a plot's thinning status can change between its own survey years, deliberately.
Yield class was inspected but excluded from the main prediction features because it is derived from forest growth information and presents a leakage or circularity concern; it is retained in the master export as an audit column and is the sole ingredient for Avenue 2's target (Methodology \S4.4.2).
Windthrow hazard class was treated as an existing forest descriptor rather than a direct observation of wind exposure, and was not included in the models until wind data were also added, since it showed a similar spatial relationship and was constant across survey years for a given plot -- confirmed as \texttt{whcl} in \texttt{MASTER\_COLUMNS}, kept as "audit/sensitivity only, never a baseline model feature" per the code's own comment.

These stand-management variables entered the main stand-based DNN and PINN comparison as the no-environment feature set.
Most environmental and terrain data (elevation, wind, CHELSA climatologies, soil) entered both the prediction models (\texttt{dnn\_env\_terrain}, \texttt{pinn\_env\_terrain}) and the later attribution analyses, so prediction and attribution were not fully separate questions for those variables.
HadUK-Grid climate variables (\texttt{tas\_mean}, \texttt{groundfrost\_mean}) are the one exception: as of this pass they are used only in the attribution-side screening (XGBoost/SHAP), because of a structural gap in the prediction pipeline's handling of cohort-specific column names, not because the variables were judged untrustworthy. \textbf{PENDING:} whether this gap has since been closed in the predictive pipeline is a Methodology/Results-chapter question, not re-checked here -- do not assume it is still open without checking the current state of \texttt{models/common/torch\_data.py} at write-up time.

\section{Terrain data}
OS Terrain 50 provided a static representation of topography at approximately 50\,m resolution.

The derived variables can be grouped by meaning: height and steepness (elevation and slope); orientation (aspect, and the northness/eastness measures derived from it); landform (terrain position index, curvature and local relief); exposure and shelter (TOPEX and related horizon measures); and local terrain heuristics (the frost-hollow flag and any investigated soil-depth proxy).

The raw aspect field is a 0--360\textdegree{} compass bearing and was not used directly as a model input, since a circular bearing has a discontinuity at 0\textdegree{}/360\textdegree{} that a model would otherwise treat as a large numeric jump between two directions that are physically adjacent.
It was instead replaced by its two continuous trigonometric components, computed from the aspect in radians ($\theta$):
\begin{align*}
\text{northness} &= \cos(\theta) \\
\text{eastness} &= \sin(\theta)
\end{align*}
so a north-facing slope (aspect $0\textdegree$) has northness $=1$, and an east-facing slope (aspect $90\textdegree$) has eastness $=1$.

Terrain is static and cannot explain changes in conditions between surveys.
TOPEX describes topographic exposure or shelter; it is not measured wind.
The frost-hollow flag is a project-derived indicator, not an observation of frost damage.
Any proposed soil-depth proxy should be labelled as exploratory unless it has strong ecological support.

\section{Wind data}
Global Wind Atlas provides raster-based, long-term wind-climate summaries.
It represents broad spatial differences in typical wind conditions and is effectively static for this analysis.
It can support spatial attribution but cannot identify conditions during a specific survey interval.

MIDAS station observations provide temporally varying weather-station measurements.
These were summarised over the intervals between LiDAR surveys, allowing investigation of mean wind, extreme wind or storm conditions during different growth periods.
They may not represent conditions at every forest plot equally, since station observations require spatial matching or interpolation.

Global Wind Atlas variables represented long-term spatial wind climate, MIDAS variables represented observed conditions over particular periods, and TOPEX represented terrain-based exposure.
These sources describe related but different aspects of wind exposure and were not treated as interchangeable.

\section{Climate, soil and other spatial data}
HadUK-Grid provided data at approximately 1\,km spatial resolution, used to investigate variables such as temperature, ground frost and rainfall.
Depending on the extraction, variables represented individual years, or were averaged across each cohort's own survey years once the multi-year extraction was in place -- confirmed in \texttt{progress\_notes.md} (2026-07-30 entry): \texttt{tas}/\texttt{groundfrost} are cohort-aware multi-year averages via \texttt{models/common/download\_haduk\_multi\_year.py}, not the earlier single-year (2021-only) extraction.
The 1\,km resolution means many nearby plots share identical values.
As a compact diagnostic example: the temperature variable (\texttt{tas}, multi-year cohort average) contained only 149 distinct values across 71,766 plots. \textbf{Confirmed directly from the notebook's own saved numeric output} (\texttt{notebooks/environmental\_data/figures/aux\_data\_resolution\_check\_results.csv}, read directly rather than trusting the notebook's prose): Moran's I $=0.9410$ ($p=0.001$, 999 permutations), between-compartment intraclass correlation $=0.8933$, Spearman vs.\ the CR residual $=0.2733$ (4survey, $p<0.001$) -- all three figures the earlier draft cited (0.941 / 0.893 / 0.273) are correct as stated. \texttt{groundfrost} (same product family) independently confirms 149 distinct values and an even higher Moran's I ($0.9862$).
It therefore represented broad spatial climate zones more strongly than fine-scale variation between neighbouring plots.

ERA5-Land was investigated as an additional reanalysis source.
Its temperature extraction produced only eight distinct values across the study plots (confirmed, \texttt{progress\_notes.md}: "ERA5: 8 distinct values, rho=0.11"), coarser than HadUK-Grid's 149.
It was not retained in the final feature sets used for modelling.

CHELSA supplied long-term (1981--2010) climatic summaries at the same nominal 1\,km grid as HadUK-Grid, but using terrain-conditioned downscaling rather than plain interpolation.
It was retained in the final feature sets used for both the attribution screening and the environment/terrain prediction models.

SoilGrids pH and the CEH soil and terrain classifications were explored as possible controls on spatial growth variation.
These are largely static variables that may describe broad site conditions but cannot establish short-term environmental effects.
The CEH categorical layers (pedotope, subsurface drainage, textural composition) each carry only a small number of classes across the plots, so the same class is assigned to many plots. \textbf{Confirmed per-layer} from the resolution-check CSV's own \texttt{n\_distinct\_values} column, not the registry's rounded "3 to 7" range: \texttt{ceh\_pedotope} 7 classes (71,727 plots sampled), \texttt{ceh\_subsurface\_drainage} 3 classes, \texttt{ceh\_textural\_composition} 4 classes -- the draft's "3 to 7" range is correct and now attributable to an exact per-layer source.

\textbf{Data-source triage criterion, answering an open question in the original plan} ("what actual method or criteria" decided whether a candidate environmental source was worth including): a \texttt{trust\_score} $\in[0,1]$, computed once per candidate source in \texttt{av1\_aux\_data\_resolution\_check.ipynb} (formula read directly from the notebook's code cell, not inferred from prose):
\[
\text{trust\_score} = \mathbb{1}[\text{data reachable}] \times \big(0.35\,(1-\text{ICC}) + 0.15\,\mathbb{1}[\text{Moran's } p<0.05] + 0.30\,\text{provenance} + 0.20\,\text{temporal match}\big)
\]
where $(1-\text{ICC})$ rewards sources with more \emph{within}-compartment variation (i.e.\ real plot-to-plot signal, not just a compartment-level flat value); \texttt{provenance} and \texttt{temporal match} are the notebook author's own explicit, hand-assigned 0--1 weights per source (e.g.\ HadUK-Grid provenance 0.75 as a station-interpolated observational product vs.\ Global Wind Atlas 0.60 as a modelled mesoscale output; disclosed in-code, not hidden). \textbf{Read critically:} this is a documented triage heuristic for deciding which sources were worth pursuing further at all (Tier 1 of the environmental-data work) -- it is distinct from, and earlier than, the two formal feature-selection procedures that actually chose the final modelling variable sets (the VIF-based screen for Avenue 1/predictive features and the R²-swap representation check for Avenue 2; both described in the Methodology chapter). Two of its four inputs are the author's own subjective judgement calls, not statistics -- appropriate for an exploratory triage step, but it should not be presented as a formal statistical selection criterion in the dissertation text.

\textbf{Critical catch, kept as a caution rather than a fact:} the same resolution-check CSV's row for the James Hutton 1:250,000 WRB soil map records \texttt{n\_plots\_sampled=0} and \texttt{n\_distinct\_values=0}, with its own saved note text still reading "SURPRISING: 0 distinct WRB soil groups found in just a 150-point sample" -- but \texttt{progress\_notes.md}'s 21 July entry, and the notebook's own hard-coded sampling logic (fixed seed, $n=150$), describe a successful earlier run of the identical cell returning \textbf{6} distinct WRB groups from a valid 150-point sample. This looks like a live external WMS call (\texttt{druid.hutton.ac.uk}) that failed silently on whatever run last regenerated the CSV, overwriting a real result with an all-zero one -- not a genuine finding that the source is uninformative. Trust neither the "0" row nor the "6 groups" narrative blindly; re-run that one notebook cell before citing a specific WRB class count. Low priority: this source is not in any final feature set used for modelling (Methodology chapter), so the discrepancy does not affect any result, only this chapter's own data-survey narrative if it were to mention a specific count.

Distances to roads, watercourses and forest boundaries were explored where relevant.
These are spatial context variables, not direct measures of growth processes.

\begin{table}[h]
\centering
\begin{tabular}{p{5.5cm}p{7.5cm}}
\hline
Limitation & Effect on interpretation \\
\hline
Complete-observation requirement & May favour stable plots that remained identifiable across all surveys \\
Four or six observations per plot & Limits estimation of detailed plot-specific trajectories \\
Unequal survey intervals & Requires annualisation and may hide within-interval changes \\
LiDAR height is not identical to biological growth & Changes may include canopy, measurement or disturbance effects \\
Limited acquisition metadata & Survey-year measurement differences cannot be fully separated \\
Coarse environmental rasters & Many plots receive the same climate value (e.g.\ 149 distinct HadUK-Grid values, 8 distinct ERA5-Land values, across 71,766 plots) \\
Static terrain and climate summaries & Describe spatial conditions more readily than temporal change \\
Station-based wind observations & Conditions at a station may differ from conditions at individual plots \\
Uncertain management history & Recorded thinning may not capture every intervention or disturbance \\
Multiple correlated environmental variables & Makes individual attribution difficult and requires controlled selection \\
Observational data & Environmental relationships indicate association, not causation \\
\hline
\end{tabular}
\end{table}

The environmental datasets provide several representations of climate, terrain and exposure, but their resolutions and time periods differ.
The attribution analysis can therefore identify variables associated with departures from expected growth, but it cannot isolate causal environmental effects.

%=======================================================================================
% VERIFICATION SUMMARY FOR THIS CHAPTER (2026-08-10, two passes) -- not for the dissertation
% itself, kept here so the next editing pass knows exactly what still needs a source. Second
% pass specifically went into the exploratory notebooks behind the environmental-data facts, per
% instruction: trust saved numeric outputs (CSVs, direct code re-runs), not notebook prose --
% narrative text in these notebooks ("SURPRISING", "REAL, not a footnote", etc.) is the author's
% own framing and was not taken at face value anywhere in this pass.
%
% CONFIRMED this session, either against code/notebook-saved-numbers, or by direct re-run:
%   - elev_percentile_95th as TARGET_COLUMN, no x1.1 correction (clean_master_data.py)
%   - six-stage cleaning funnel order and exact thresholds (MIN_PLYR=1850, Age in [0,200],
%     height in (0,60]) (clean_master_data.py)
%   - ALL cleaning-funnel row/plot counts, by direct re-run of clean_master_data.py against the
%     raw GPKG (2026-08-10) -- 4survey and 6survey both match the table exactly; additionally
%     established that stages 5-6 remove 0 rows in EITHER cohort, and stage 4 removes rows only
%     in 4survey (6survey loses nothing from stage 3 onward) -- a sharper fact than the table
%     alone states
%   - Thin/time_since_thinning/thinning_status derivation logic (clean_master_data.py)
%   - whcl kept audit-only, never a baseline feature (clean_master_data.py comment)
%   - HadUK-Grid (tas): 149 distinct values, Moran's I=0.9410, ICC=0.8933, Spearman=0.2733 --
%     all read directly from aux_data_resolution_check_results.csv (the notebook's saved numeric
%     output), not its prose; groundfrost independently confirms 149 distinct values
%   - ERA5-Land 8 distinct values (same CSV)
%   - ceh_pedotope 7 classes, ceh_subsurface_drainage 3 classes, ceh_textural_composition 4
%     classes -- each confirmed individually from the CSV's own n_distinct_values column
%   - Global Wind Atlas: 1,816 distinct values (CSV); native cell size recomputed directly from
%     the cached raster (rasterio) = 154.9m E-W x 278.3m N-S at Aberfoyle's latitude -- this
%     recomputation is what resolved a disagreement between two notebooks (one said 250m N-S,
%     the correct one said 278m)
%   - trust_score formula (data-source triage criterion) extracted directly from the notebook's
%     code cell -- answers the original plan's own open question about selection criteria; noted
%     as a distinct, earlier, partly-subjective step from the later VIF/R2-swap feature selection
%
% STILL OPEN / FLAGGED AS A CAUTION, NOT A FACT:
%   - James Hutton WRB soil map: saved CSV row (n=0 distinct groups) directly contradicts
%     progress_notes.md's narrative (6 groups from a valid 150-point run) -- diagnosed as a
%     likely silent live-API failure on whatever run last regenerated the CSV, not a real
%     zero-information finding. Not resolved (would need re-running that one notebook cell
%     against the live WMS); low priority since this source isn't in any final feature set.
%   - LAI column: no plausibility filter applied in current cleaning code; whether an
%     earlier-flagged anomaly in this column is still present in the current cleaned cohorts
%     remains unchecked (not investigated in the notebooks reviewed this pass either).
%=======================================================================================
